% Guide d'utilisation — Thème Proxmox Cyberpunk V28 (portable)
% Compilation : lualatex guide_theme_proxmox.tex   (depuis ce dossier)
% Polices : Lato (paquet Debian fonts-lato), DejaVu Sans Mono, Orbitron (fournie ici, SIL OFL 1.1)
\documentclass[9pt]{extarticle}
\usepackage[a4paper,margin=11mm,top=9mm,bottom=10mm]{geometry}
\usepackage{fontspec}
\usepackage[french]{babel}
\usepackage{xcolor,graphicx,tikz,multicol,enumitem,tabularx,array}
\usepackage[most]{tcolorbox}
\usepackage{microtype}
\usepackage{amssymb}

\setmainfont{Lato}[UprightFont=*-Regular,BoldFont=*-Bold,ItalicFont=*-Italic,BoldItalicFont=*-BoldItalic]
\setmonofont{DejaVu Sans Mono}[Scale=0.84]
\newfontfamily\orbitron{Orbitron-900.ttf}[Path=./]
\newfontfamily\orbitronb{Orbitron-700.ttf}[Path=./]
\newfontfamily\latoblack{Lato-Black}

\definecolor{nuit}{HTML}{0A0E1A}
\definecolor{nuit2}{HTML}{14233F}
\definecolor{cyan}{HTML}{00B8D9}
\definecolor{cyanvif}{HTML}{00E5FF}
\definecolor{vert}{HTML}{00A86B}
\definecolor{rouge}{HTML}{CC2244}
\definecolor{encre}{HTML}{1B2433}
\definecolor{gris}{HTML}{5B6676}
\definecolor{fondcode}{HTML}{0F1A2E}

\graphicspath{{../captures/}}
\pagestyle{empty}
\setlength{\parindent}{0pt}
\setlength{\parskip}{2pt}
\setlength{\columnsep}{7mm}
\color{encre}
\setlist{nosep,leftmargin=11pt,itemsep=1pt,label={\color{cyan}\scriptsize$\blacktriangleright$}}

% Titre de rubrique
\newcommand{\rubrique}[1]{%
  \par\vspace{5pt}{\orbitronb\fontsize{9.6}{11}\selectfont\color{nuit2}\MakeUppercase{#1}}\par
  \vspace{-2pt}{\color{cyan}\rule{\linewidth}{0.9pt}}\par\vspace{2pt}}

% Bloc de commandes (fond sombre comme le thème)
\newtcolorbox{commande}{enhanced,colback=fondcode,colframe=cyan!70!black,boxrule=0.5pt,arc=2pt,
  left=4pt,right=4pt,top=2.5pt,bottom=2.5pt,fontupper=\ttfamily\color{cyanvif!35!white},before skip=3pt,after skip=4pt}
\newcommand{\cmt}[1]{{\color{gris!80!white}\##1}}
\newcommand{\opt}[1]{{\ttfamily\bfseries\color{nuit2}#1}}
\newcommand{\optt}[1]{{\ttfamily\bfseries\color{nuit2}\fontsize{7.6}{9}\selectfont #1}}
\newcommand{\chemin}[1]{{\ttfamily\small #1}}

\begin{document}

% ───────────── Bandeau ─────────────
\begin{tikzpicture}[remember picture,overlay]
  \fill[nuit] (current page.north west) rectangle ([yshift=-31mm]current page.north east);
  \shade[left color=nuit, right color=nuit2] ([yshift=-31mm]current page.north west) rectangle ([yshift=-26mm]current page.north east);
  \fill[cyanvif] ([yshift=-31mm]current page.north west) rectangle ([yshift=-31.8mm]current page.north east);
  \foreach \x/\y/\r in {0.12/0.35/0.5, 0.31/0.7/0.35, 0.55/0.25/0.45, 0.72/0.62/0.3, 0.88/0.4/0.5, 0.95/0.75/0.3}
    \fill[cyanvif,opacity=0.55] ([xshift=\x*\paperwidth,yshift=-\y*26mm]current page.north west) circle (\r pt);
\end{tikzpicture}
\vspace*{-1mm}
{\orbitron\fontsize{22}{24}\selectfont\color{cyanvif}THÈME PROXMOX CYBERPUNK}\par
\vspace{1.5mm}
{\fontsize{11}{13}\selectfont\color{white}\textbf{Guide d'installation pour collègues} \;—\; version portable V28 \;—\; Proxmox VE 8 et 9}\par
\vspace{0.8mm}
{\small\color{cyanvif!60!white}Damien SCONTRINO — BTS SIO, lycée Louis Armand (Nogent-sur-Marne) — usage pédagogique réservé au réseau CERTA}\par
\vspace{7.5mm}

\begin{multicols}{2}

\rubrique{À quoi sert ce script ?}
Il habille l'interface web de Proxmox VE en \textbf{bleu néon} : fond sombre animé, icônes colorées
selon l'état ({\color{rouge}\textbf{VM arrêtée}}, VM en marche, \textcolor{violet}{\textbf{modèle}}), menus, fenêtres et
graphiques assortis, et un \textbf{badge} dans la barre du haut (nom de l'établissement, nom du nœud…).
Trois \textbf{options pédagogiques}, désactivées par défaut : message d'accueil après la connexion,
« clone intégral » imposé aux étudiants, retrait de la fenêtre « No valid subscription ».\par
\textbf{Aucun téléchargement, aucun service redémarré, retour arrière en une commande.}

\rubrique{Prérequis}
\begin{itemize}
  \item Proxmox VE \textbf{8.x} (Debian 12) ou \textbf{9.x} (Debian 13), nœud seul ou cluster, n'importe quel nom de nœud.
  \item Un accès \textbf{root} (console, shell web ou SSH). Rien à installer : bash, perl et coreutils suffisent.
  \item En cluster : la clé SSH root entre nœuds (déjà en place dans tout cluster Proxmox).
\end{itemize}

\rubrique{Installation}
Copier le script sur le serveur, puis en root :
\begin{commande}
chmod +x install\_theme\_proxmox.sh\\
./install\_theme\_proxmox.sh \cmt{ menu interactif}
\end{commande}
Le menu propose : \textbf{1} installer / mettre à jour, \textbf{2} désinstaller, \textbf{3} état,
\textbf{4} simulation (ne modifie rien). Il pose les questions (badge, options), affiche un récapitulatif
et demande confirmation. Ensuite : \textbf{Ctrl+F5} dans le navigateur.
Installation directe, sans question :
\begin{commande}
./install\_theme\_proxmox.sh --auto --badge "\%NODE\% - Lycée X"
\end{commande}
Relancer le script avec d'autres options \textbf{met à jour} l'installation : jamais de doublon.


\columnbreak

\rubrique{Ce qui est modifié (et rien d'autre)}
\begin{itemize}
  \item \chemin{index.html.tpl} : 1 ligne \texttt{<link>} (+1 \texttt{<script>} si option), repérées par la marque \texttt{theme-proxmox-cyberpunk} ;
  \item une feuille de style dédiée (police Orbitron intégrée) et, si option, un petit script ;
  \item \chemin{proxmoxlib.js} \emph{uniquement} avec \opt{--sans-abonnement} ;
  \item un crochet APT et le dossier \chemin{/opt/theme-proxmox-cyberpunk/}.
\end{itemize}
\chemin{ext6-pve.css} et \chemin{pvemanagerlib.js} restent intacts. Sauvegarde avant chaque modification dans
\chemin{/root/sauvegardes-theme-proxmox/}, journal \chemin{/var/log/theme-proxmox-cyberpunk.log}.

\rubrique{Désinstallation}
\begin{commande}
/opt/theme-proxmox-cyberpunk/install\_theme\_proxmox.sh --desinstaller
\end{commande}
Ajouter \opt{--cluster} pour tous les nœuds. Le script retire ses lignes et ses fichiers, puis \textbf{vérifie}
que chaque fichier de Proxmox est redevenu identique à celui du paquet Debian (empreinte md5 de dpkg).
Il ne touche pas à ce qu'il n'a pas modifié lui-même.

\rubrique{Mises à jour de Proxmox}
Une mise à jour de \texttt{pve-manager} réécrit \chemin{index.html.tpl} : le \textbf{crochet APT réapplique
le thème automatiquement} à la fin d'\texttt{apt} (« Thème Proxmox cyberpunk : réappliqué après la mise à
jour »). Le passage de PVE 8 à 9 est suivi. Sur une version future non testée (10.x), le thème n'est pas
réappliqué : l'interface d'origine reste (option \opt{--forcer} pour essayer).

\rubrique{Cluster}
L'interface est servie par \textbf{chaque} nœud : le menu propose d'installer partout (ou \opt{--cluster}).
Un nœud hors ligne est signalé et ignoré ; relancer plus tard sur ce nœud.

\rubrique{Dépannage}
\begin{itemize}
  \item \textbf{Thème invisible} : Ctrl+F5 (cache du navigateur), puis \opt{--etat}.
  \item \textbf{Ancien thème encore présent} (couleurs en double, ancien pop-up) : \opt{--nettoyer-ancien}.
  \item \textbf{« </head> trouvée 0 fois »} : modèle modifié par un autre outil, le script refuse d'y toucher.
  \item \textbf{Doute ou interface étrange} : \opt{--desinstaller} ; en dernier recours\\
        \texttt{apt install --reinstall pve-manager proxmox-widget-toolkit}.
  \item \textbf{Message d'accueil} : une fois par session de navigateur ; il revient si son texte change.
  \item \textbf{Nœud en échec en cluster} : copier le script sur ce nœud et le lancer localement.
\end{itemize}

\vspace{4pt}
\begin{tcolorbox}[enhanced,colback=cyan!6,colframe=cyan!60!black,boxrule=0.4pt,arc=2pt,left=4pt,right=4pt,top=2pt,bottom=2pt,fontupper=\scriptsize]
\textbf{Testé} : installation, relance, mise à jour réelle de \texttt{pve-manager} (crochet APT) et désinstallation
sur PVE~8.4 (nœud seul, fichiers d'origine retrouvés à l'identique, sha256) ; mêmes étapes sur PVE~9.2 en bac
à sable avec les fichiers du paquet. Captures d'écran réelles page suivante.
\end{tcolorbox}

\end{multicols}

\vfill
{\scriptsize\color{gris}Licence : usage pédagogique réservé aux enseignants du réseau CERTA — © 2025-2026 Damien SCONTRINO
(damien.scontrino@ac-creteil.fr). Police Orbitron © 2018 The Orbitron Project Authors, SIL Open Font License 1.1 (\texttt{OFL-Orbitron.txt}).}

% ───────────── Page 2 : captures ─────────────
\newpage
\begin{tikzpicture}[remember picture,overlay]
  \fill[nuit] (current page.north west) rectangle ([yshift=-15mm]current page.north east);
  \fill[cyanvif] ([yshift=-15mm]current page.north west) rectangle ([yshift=-15.6mm]current page.north east);
\end{tikzpicture}
\vspace*{-0.5mm}
{\orbitron\fontsize{14}{16}\selectfont\color{cyanvif}CAPTURES D'ÉCRAN}\hfill{\small\color{white}Proxmox VE 8.4, nœud de test, badge « \%NODE\% - Lycée Exemple »}\par
\vspace{8mm}

\newcommand{\cadre}[1]{%
  \begin{tcolorbox}[enhanced,colback=nuit,colframe=cyan!70!black,boxrule=0.6pt,arc=2pt,left=0pt,right=0pt,top=0pt,bottom=0pt,
    drop fuzzy shadow=cyan!40!black]
    \includegraphics[width=\linewidth]{#1}
  \end{tcolorbox}}

\cadre{3_interface.png}\vspace{-3pt}
{\small\textbf{Vue d'un nœud.} Arborescence : VM arrêtée en rouge, modèles en violet, conteneur, réseau et stockages ;
logo néon, badge (ici \texttt{\%NODE\% - Lycée Exemple}), menus, grilles et barre des tâches au même style.}\par
\vspace{4mm}
\begin{minipage}[t]{0.488\linewidth}
  \cadre{1_connexion_zoom.png}\vspace{-3pt}
  {\small\textbf{Page de connexion.} Le thème s'applique dès l'ouverture de l'interface.}
\end{minipage}\hfill
\begin{minipage}[t]{0.488\linewidth}
  \cadre{2_message_accueil_zoom.png}\vspace{-3pt}
  {\small\textbf{Message d'accueil} (option \opt{--message}, fichier d'exemple fourni), affiché après la connexion.}
\end{minipage}\par
\vspace{5mm}

\begin{tcolorbox}[enhanced,colback=cyan!6,colframe=cyan!60!black,boxrule=0.4pt,arc=2pt,left=5pt,right=5pt,top=3pt,bottom=3pt]
\small\textbf{À propos des captures.} Prises avec un navigateur sans écran sur un nœud de test, avec un compte
temporaire en lecture seule supprimé ensuite (aucun identifiant réel). Les pools et stockages nommés d'après
des comptes d'étudiants ont été retirés de l'affichage ; le reste est l'interface réelle avec le thème.
\end{tcolorbox}

\vspace{4mm}
\rubrique{Fichiers fournis}
\begin{tabularx}{\linewidth}{@{}>{\ttfamily}p{55mm}X@{}}
install\_theme\_proxmox.sh & le script, autonome (police intégrée) : c'est le seul fichier indispensable\\
exemple\_message\_accueil.html & exemple pour \opt{--message}, à adapter (HTML simple ou texte brut)\\
OFL-Orbitron.txt & licence de la police Orbitron intégrée au thème\\
guide\_theme\_proxmox.pdf & ce guide (source LaTeX dans \texttt{source\_guide/})\\
captures/ & les captures en taille réelle (dont la fiche d'une VM, \texttt{4\_vm.png})\\
\end{tabularx}

\end{document}
